\subsection{Business models}

\begin{qfdefinition}[Business models: the value proposition, customers, channels, key resources and partners, revenue model and cost structure, the distinction from strategy, and common types including platforms, subscriptions and asset-light models]\label{def:business-model-features-types}
\qffield{Statement} A business model describes how a company creates value for chosen customers, delivers it and captures part of it as profit. Osterwalder, Pigneur and Tucci synthesise the literature into nine building blocks under four pillars \citep[pp.~17--18 and Table~3]{clarifying-business-models-2005}:
\begin{itemize}
\item \emph{product}: the value proposition;
\item \emph{customer interface}: target customers, distribution channels and relationships;
\item \emph{infrastructure}: the value configuration of activities, core competencies and the partner network;
\item \emph{financial aspects}: the cost structure and the revenue model.
\end{itemize}
They distinguish three levels of the concept: the general definition, taxonomies of model types, and descriptions of particular companies (pp.~8--10). A business model is not a strategy. The model describes how the pieces of a business fit together. Strategy adds how the company will compete against rivals and how it will execute (p.~13, citing Magretta).

\emph{Types.} Common patterns, each with its own value drivers (Interpretation):
\begin{itemize}
\item \emph{Product sale through channels}: revenue per unit sold, with margins shared with distributors.
\item \emph{Subscription and recurring revenue}: payment for continued access. Value depends on retention and the cost of acquiring customers.
\item \emph{Razor and blades}: a base product sold cheaply and profit taken on consumables or services.
\item \emph{Asset-light and franchise}: brand and system owned, assets owned by partners.
\item \emph{Platforms and two-sided markets}: the company enables interactions between groups of users and must get both sides on board. A market is two-sided when the volume of transactions depends on how the price is split between the sides, not only on its total. Examples include video-game consoles, operating systems, advertising-supported media and payment cards \citep[pp.~1--2]{two-sided-markets-2005}.
\end{itemize}

Damodaran's treatment of Uber in 2014 shows the model feeding a valuation. Uber was a matchmaker, not a taxi company. It took a slice, about 20 percent, of each fare, owned no cars, and reinvested in technology and acquisitions. That made capital intensity low and margins dependent on the take rate, which competition could cut \citep[pp.~83--85]{narrative-and-numbers-2017}.

\qffield{Practice} \emph{Original.} A ride platform processes gross bookings of 1,000 at a 20 percent take rate, earning revenue of 200. If competition forces the rate to 15 percent, bookings must rise by a third to hold revenue; a 20 percent rise gives $1{,}200\times0.15=180$. A valuation should therefore model bookings, take rate and capital intensity separately, rather than revenue alone.

\emph{Harbour Tools} (\hyperref[ex:harbour-tools-statements]{case}) sells power tools and batteries through distributors and retailers. Applying the building blocks:
\begin{itemize}
\item \emph{Value proposition}: reliable tools for trade and home users.
\item \emph{Customers and channels}: two segments, reached through distributors and retailers with whom margin is shared, and limited direct relationships.
\item \emph{Infrastructure}: owned plants and capitalised product development.
\item \emph{Revenue model}: one-off product sales priced per unit, with battery platforms encouraging repeat purchases within the brand.
\item \emph{Cost structure}: materials, labour and fixed plant costs, giving operating leverage.
\end{itemize}
The value drivers this points to are volume, price, gross margin, capital and working-capital intensity, and the reinvestment needed for growth. These are the drivers in the \hyperref[ex:harbour-pro-forma-model]{pro forma model}. A shift to tool rental by subscription would change the revenue model, raise capital intensity and make customer retention a driver.

\qffield{Conditions} Business-model taxonomies are descriptive. The same label covers firms with very different economics, and a company can run several models at once, in which case each segment needs its own description. The two-sided-market definition is economic, so it applies only where the price split, not only the price level, affects volume.

\qffield{Pitfalls} Treating a fashionable label, such as ``platform'' or ``asset-light'', as evidence of value. Describing the model without linking it to the forecast drivers it implies. Confusing the business model with strategy, and so missing how competitors could copy the model. Valuing one side of a platform as if the other side did not respond to its prices.

\qffield{See also} \hyperref[meth:industry-analysis]{industry analysis}; \hyperref[meth:company-position-analysis]{company position}; \hyperref[ex:harbour-pro-forma-model]{Harbour pro forma model}.

\end{qfdefinition}
